\subsection{The Aggregate Welfare Function}
\label{subsec:aggregate-welfare}

%%%%%%%%%%%%%%%%%%%%%%%%%%%%%%%%%%%%%%%%%%%%%%%%%%%%%%%%%%%%%%%%%%%%%%%%%%%%%%%
% 10.8: Putting it all together - the complete EBF welfare function
%%%%%%%%%%%%%%%%%%%%%%%%%%%%%%%%%%%%%%%%%%%%%%%%%%%%%%%%%%%%%%%%%%%%%%%%%%%%%%%

This section integrates all components developed in Sections~\ref{subsec:three-categories}--\ref{subsec:context-modulation} 
into the complete EBF welfare function.

\subsubsection{The Integration Challenge}

\paragraph{What We Have Established.}
The previous sections developed:
\begin{enumerate}
\item Three utility categories: INU, KNU, IDN (Section~\ref{subsec:three-categories})
\item Six FEPSDE dimensions: F, E, P, S, D, Eco (Section~\ref{subsec:fepsde-dimensions})
\item 144-component structure: $3 \times 6 \times 4 \times 2$ (Section~\ref{subsec:144-structure})
\item Uniform calculation logic: $U^i_{d,t} = \delta^t \cdot [G - \lambda \cdot P]$ (Section~\ref{subsec:calculation-logic})
\item Reference point formation: Category-specific $r^i_d$ (Section~\ref{subsec:reference-points})
\item Inter-category complementarities: $\gamma_{ij}$ (Section~\ref{subsec:inter-category-complementarities})
\item Context modulation: $\Gamma$-matrix (Section~\ref{subsec:context-modulation})
\end{enumerate}

\paragraph{The Integration Task.}
Now we must combine these elements into a single, coherent welfare function 
that can be used for:
\begin{itemize}
\item Individual welfare assessment
\item Policy evaluation
\item Behavioral change prediction
\item Intervention design
\end{itemize}

%%%%%%%%%%%%%%%%%%%%%%%%%%%%%%%%%%%%%%%%%%%%%%%%%%%%%%%%%%%%%%%%%%%%%%%%%%%%%%%
\subsubsection{Level 1: Component-Level Utility}
%%%%%%%%%%%%%%%%%%%%%%%%%%%%%%%%%%%%%%%%%%%%%%%%%%%%%%%%%%%%%%%%%%%%%%%%%%%%%%%

\paragraph{The Atomic Unit.}
The foundation is the individual component utility:

\begin{equation}
\boxed{
U^i_{d,t}(\Psi) = \left(\delta^i_d(\Psi)\right)^t \cdot \left[ G^i_{d,t} - \lambda^i_d(\Psi) \cdot P^i_{d,t} \right]
}
\label{eq:component-utility}
\end{equation}

where:
\begin{align*}
G^i_{d,t} &= \max\left(0, x^i_{d,t} - r^i_d(\Psi)\right) \\
P^i_{d,t} &= \max\left(0, r^i_d(\Psi) - x^i_{d,t}\right)
\end{align*}

This yields 144 component utilities, one for each combination of:
\begin{itemize}
\item Category $i \in \{INU, KNU, IDN\}$
\item Dimension $d \in \{F, E, P, S, D, Eco\}$
\item Time $t \in \{0, 1, 2, 3\}$
\item Valence (embedded in $G$ vs. $P$)
\end{itemize}

%%%%%%%%%%%%%%%%%%%%%%%%%%%%%%%%%%%%%%%%%%%%%%%%%%%%%%%%%%%%%%%%%%%%%%%%%%%%%%%
\subsubsection{Level 2: Category-Level Utility}
%%%%%%%%%%%%%%%%%%%%%%%%%%%%%%%%%%%%%%%%%%%%%%%%%%%%%%%%%%%%%%%%%%%%%%%%%%%%%%%

\paragraph{Aggregating Within Categories.}
For each category, we aggregate across dimensions and time:

\begin{equation}
U^i(\Psi) = \sum_{d \in FEPSDE} \sum_{t=0}^{3} \nu^i_d \cdot U^i_{d,t}(\Psi)
\label{eq:category-utility}
\end{equation}

where $\nu^i_d$ is the dimension weight within category $i$.

\paragraph{Dimension Weights.}
The weights $\nu^i_d$ may vary by category:

\begin{table}[htbp]
\centering
\caption{Dimension Weights by Category (Illustrative)}
\label{tab:dimension-weights}
\renewcommand{\arraystretch}{1.3}
\begin{tabular}{@{}lcccccc@{}}
\toprule
& $\nu_F$ & $\nu_E$ & $\nu_P$ & $\nu_S$ & $\nu_D$ & $\nu_{Eco}$ \\
\midrule
INU & 0.25 & 0.20 & 0.20 & 0.15 & 0.10 & 0.10 \\
KNU & 0.20 & 0.15 & 0.15 & 0.25 & 0.10 & 0.15 \\
IDN & 0.15 & 0.30 & 0.15 & 0.25 & 0.05 & 0.10 \\
\bottomrule
\end{tabular}
\end{table}

Note: These weights are person-specific and context-dependent. The table 
shows illustrative population averages.

\paragraph{Result.}
This yields three category-level utilities:
\begin{align}
U^{INU}(\Psi) &= \text{Total individual utility} \\
U^{KNU}(\Psi) &= \text{Total collective utility} \\
U^{IDN}(\Psi) &= \text{Total identity utility}
\end{align}

%%%%%%%%%%%%%%%%%%%%%%%%%%%%%%%%%%%%%%%%%%%%%%%%%%%%%%%%%%%%%%%%%%%%%%%%%%%%%%%
\subsubsection{Level 3: Total Welfare with Complementarities}
%%%%%%%%%%%%%%%%%%%%%%%%%%%%%%%%%%%%%%%%%%%%%%%%%%%%%%%%%%%%%%%%%%%%%%%%%%%%%%%

\paragraph{The Naive Aggregation (Insufficient).}
A first attempt might simply weight the categories:

\begin{equation}
Q_{naive} = \omega_{INU} \cdot U^{INU} + \omega_{KNU} \cdot U^{KNU} + \omega_{IDN} \cdot U^{IDN}
\end{equation}

But this ignores inter-category interactions.

\paragraph{The Full Welfare Function.}
Including complementarities:

\begin{equation}
\boxed{
Q(\Psi) = \sum_{i} \omega_i \cdot U^i(\Psi) + \sum_{i < j} \gamma_{ij}(\Psi) \cdot U^i(\Psi) \cdot U^j(\Psi)
}
\label{eq:total-welfare}
\end{equation}

Expanding:
\begin{align}
Q(\Psi) = \quad & \omega_{INU} \cdot U^{INU}(\Psi) \nonumber \\
+ \quad & \omega_{KNU} \cdot U^{KNU}(\Psi) \nonumber \\
+ \quad & \omega_{IDN} \cdot U^{IDN}(\Psi) \nonumber \\
+ \quad & \gamma_{INU,KNU}(\Psi) \cdot U^{INU}(\Psi) \cdot U^{KNU}(\Psi) \nonumber \\
+ \quad & \gamma_{INU,IDN}(\Psi) \cdot U^{INU}(\Psi) \cdot U^{IDN}(\Psi) \nonumber \\
+ \quad & \gamma_{KNU,IDN}(\Psi) \cdot U^{KNU}(\Psi) \cdot U^{IDN}(\Psi)
\label{eq:welfare-expanded}
\end{align}

%%%%%%%%%%%%%%%%%%%%%%%%%%%%%%%%%%%%%%%%%%%%%%%%%%%%%%%%%%%%%%%%%%%%%%%%%%%%%%%
\subsubsection{The Complete EBF Welfare Equation}
%%%%%%%%%%%%%%%%%%%%%%%%%%%%%%%%%%%%%%%%%%%%%%%%%%%%%%%%%%%%%%%%%%%%%%%%%%%%%%%

\paragraph{Fully Expanded Form.}
Substituting all components:

\begin{tcolorbox}[colback=blue!5!white, colframe=blue!75!black, title=The EBF Welfare Function]
\begin{equation}
Q(\mathbf{x}, \Psi) = \sum_{i} \omega_i \cdot \left[ \sum_{d} \sum_{t} \nu^i_d \cdot \delta^{i,t}_d(\Psi) \cdot \left( G^i_{d,t} - \lambda^i_d(\Psi) \cdot P^i_{d,t} \right) \right] + \sum_{i<j} \gamma_{ij}(\Psi) \cdot U^i \cdot U^j
\label{eq:bcm-welfare-complete}
\end{equation}

where:
\begin{itemize}
\item $\mathbf{x}$: outcome vector (144 components)
\item $\Psi$: context vector (8 dimensions)
\item $\omega_i$: category weights (3 parameters)
\item $\nu^i_d$: dimension weights within categories (18 parameters)
\item $\delta^i_d(\Psi)$: discount factors (18 base + context modulation)
\item $\lambda^i_d(\Psi)$: loss aversion (18 base + context modulation)
\item $r^i_d(\Psi)$: reference points (18 formation rules)
\item $\gamma_{ij}(\Psi)$: complementarities (3 base + context modulation)
\end{itemize}
\end{tcolorbox}

\paragraph{Level-Indexed KNU.}
An important refinement: $U^{KNU}$ is itself aggregated across levels
(see Section~\ref{subsec:three-categories} and Appendix~AAA):

\begin{equation}
U^{KNU}(\Psi) = \sum_{L=1}^{\infty} \alpha^L(\Psi) \cdot U^{KNU,L}
\label{eq:knu-level-aggregation}
\end{equation}

where $\alpha^L(\Psi)$ is the \textbf{level salience function} from CORE-WHO.
This means ``collective utility'' depends on \emph{which collective} is psychologically
salient---family ($L=2$), organization ($L=4$), nation ($L=5$), etc.

\textit{Cross-Reference:} The 9-level hierarchy and salience functions are
formalized in Appendix~AAA (CORE-WHO: Who has utility?).

%%%%%%%%%%%%%%%%%%%%%%%%%%%%%%%%%%%%%%%%%%%%%%%%%%%%%%%%%%%%%%%%%%%%%%%%%%%%%%%
\subsubsection{Properties of the Welfare Function}
%%%%%%%%%%%%%%%%%%%%%%%%%%%%%%%%%%%%%%%%%%%%%%%%%%%%%%%%%%%%%%%%%%%%%%%%%%%%%%%

\paragraph{Property 1: Reference-Dependence.}
The welfare function is reference-dependent, not outcome-dependent:

\begin{equation}
Q(\mathbf{x}) \neq Q(\mathbf{x}') \quad \text{even if } \sum_i x_i = \sum_i x'_i
\end{equation}

Two outcomes with the same objective sum can yield different welfare 
if they differ in their relationship to reference points.

\paragraph{Property 2: Loss Aversion.}
Losses hurt more than equivalent gains help:

\begin{equation}
\left| \frac{\partial Q}{\partial P^i_d} \right| = \lambda^i_d \cdot \left| \frac{\partial Q}{\partial G^i_d} \right| > \left| \frac{\partial Q}{\partial G^i_d} \right|
\end{equation}

Since $\lambda > 1$, reducing a Pain by one unit increases welfare more 
than increasing a Gain by one unit.

\paragraph{Property 3: Diminishing Sensitivity.}
Due to prospect theory's value function shape (not explicitly modeled here 
but compatible with the framework), sensitivity to changes decreases 
as outcomes move further from the reference point.

\paragraph{Property 4: Context-Dependence.}
All parameters respond to context:

\begin{equation}
\frac{\partial Q}{\partial \Psi_j} \neq 0 \quad \text{for most } j
\end{equation}

The same outcome produces different welfare in different contexts.

\paragraph{Property 5: Non-Separability.}
Due to complementarities, the welfare contribution of one category 
depends on the levels of other categories:

\begin{equation}
\frac{\partial^2 Q}{\partial U^i \partial U^j} = \gamma_{ij} \neq 0
\end{equation}

Categories cannot be optimized independently.

%%%%%%%%%%%%%%%%%%%%%%%%%%%%%%%%%%%%%%%%%%%%%%%%%%%%%%%%%%%%%%%%%%%%%%%%%%%%%%%
\subsubsection{The Coherence-Welfare Relationship}
%%%%%%%%%%%%%%%%%%%%%%%%%%%%%%%%%%%%%%%%%%%%%%%%%%%%%%%%%%%%%%%%%%%%%%%%%%%%%%%

\paragraph{Connecting K and Q.}
The Complementarity-Context Framework's coherence measure $K$ relates to 
welfare $Q$:

\begin{equation}
\frac{\partial Q}{\partial K} > 0
\end{equation}

Higher coherence produces higher welfare, holding outcomes constant.

\paragraph{The K-Q-$\Psi$ Triad.}
The three core constructs form a causal chain:

\begin{equation}
\Psi \longrightarrow K \longrightarrow Q
\end{equation}

\begin{itemize}
\item Context ($\Psi$) shapes the possibility of coherence
\item Coherence ($K$) enables welfare realization
\item Welfare ($Q$) is the ultimate outcome
\end{itemize}

\paragraph{Formal Relationship.}
\begin{equation}
Q = f(K, \mathbf{x}, \Psi) = K \cdot \bar{Q}(\mathbf{x}) + (1-K) \cdot \underline{Q}(\mathbf{x}) + \epsilon(\Psi)
\end{equation}

where:
\begin{itemize}
\item $\bar{Q}$: welfare under full coherence
\item $\underline{Q}$: welfare under full incoherence
\item $\epsilon(\Psi)$: context-specific adjustment
\end{itemize}

High coherence ($K \rightarrow 1$) allows outcomes to translate efficiently 
into welfare. Low coherence ($K \rightarrow 0$) means even good outcomes 
produce suboptimal welfare due to internal conflicts.

%%%%%%%%%%%%%%%%%%%%%%%%%%%%%%%%%%%%%%%%%%%%%%%%%%%%%%%%%%%%%%%%%%%%%%%%%%%%%%%
\subsubsection{Practical Application: Welfare Assessment}
%%%%%%%%%%%%%%%%%%%%%%%%%%%%%%%%%%%%%%%%%%%%%%%%%%%%%%%%%%%%%%%%%%%%%%%%%%%%%%%

\paragraph{The Assessment Protocol.}
To assess welfare for an individual or population:

\begin{enumerate}
\item \textbf{Measure Context:} Estimate $\Psi$ vector (8 dimensions)
\item \textbf{Identify Category Weights:} Estimate $\omega_{INU}$, $\omega_{KNU}$, $\omega_{IDN}$
\item \textbf{Determine Reference Points:} Calculate $r^i_d(\Psi)$ for each category-dimension
\item \textbf{Observe Outcomes:} Measure $x^i_{d,t}$ for all 144 components
\item \textbf{Calculate Components:} Apply uniform calculation logic to get $U^i_{d,t}$
\item \textbf{Aggregate Within Categories:} Sum to get $U^{INU}$, $U^{KNU}$, $U^{IDN}$
\item \textbf{Apply Complementarities:} Combine with $\gamma_{ij}$ to get total $Q$
\end{enumerate}

\paragraph{Simplified Assessment (Practical Approximation).}
For practical applications, a reduced form may suffice:

\begin{equation}
Q \approx \omega_{INU} \cdot U^{INU} + \omega_{KNU} \cdot U^{KNU} + \omega_{IDN} \cdot U^{IDN} + \gamma \cdot K^{utility}
\end{equation}

where $K^{utility}$ is the utility coherence index from Section~\ref{subsec:inter-category-complementarities}.

%%%%%%%%%%%%%%%%%%%%%%%%%%%%%%%%%%%%%%%%%%%%%%%%%%%%%%%%%%%%%%%%%%%%%%%%%%%%%%%
\subsubsection{Policy Implications}
%%%%%%%%%%%%%%%%%%%%%%%%%%%%%%%%%%%%%%%%%%%%%%%%%%%%%%%%%%%%%%%%%%%%%%%%%%%%%%%

\paragraph{Implication 1: Multi-Dimensional Optimization.}
Policy should optimize across all six FEPSDE dimensions, not just GDP ($U^F$):

\begin{equation}
\max_{\text{policy}} Q = \max_{\text{policy}} \left[ \sum_d \sum_i \omega_i \nu^i_d U^i_d + \text{complementarities} \right]
\end{equation}

\paragraph{Implication 2: Category Balance.}
Policies that boost INU while harming KNU or IDN may reduce total welfare:

\begin{equation}
\Delta Q = \omega_{INU} \Delta U^{INU} + \gamma_{INU,KNU} U^{KNU} \Delta U^{INU} + \gamma_{INU,IDN} U^{IDN} \Delta U^{INU}
\end{equation}

If $\gamma < 0$, INU gains that harm other categories reduce total welfare.

\paragraph{Implication 3: Context Matters.}
The same policy produces different welfare effects in different contexts:

\begin{equation}
\frac{\partial \Delta Q}{\partial \Psi} \neq 0
\end{equation}

Policies must be calibrated to context, not applied uniformly.

\paragraph{Implication 4: Reference Point Management.}
Policy can target reference points, not just outcomes:

\begin{equation}
\Delta Q = \frac{\partial Q}{\partial x} \Delta x + \frac{\partial Q}{\partial r} \Delta r
\end{equation}

Managing expectations ($\Delta r$) can be as effective as changing outcomes ($\Delta x$).

\paragraph{Implication 5: Sequencing Principle.}
Given the complementarity structure, policy sequencing matters:

\begin{tcolorbox}[colback=yellow!10!white, colframe=yellow!50!black, title=The Sequencing Principle]
\begin{enumerate}
\item \textbf{First:} Establish context ($\Psi$)---institutions, trust, stability
\item \textbf{Second:} Build complementarities ($\gamma > 0$)---align categories
\item \textbf{Third:} Improve outcomes ($x$)---now they translate efficiently to welfare
\end{enumerate}

Improving outcomes without adequate context and complementarity structure 
yields suboptimal welfare gains.
\end{tcolorbox}

%%%%%%%%%%%%%%%%%%%%%%%%%%%%%%%%%%%%%%%%%%%%%%%%%%%%%%%%%%%%%%%%%%%%%%%%%%%%%%%
\subsubsection{Connecting to Behavioral Change (Preview)}
%%%%%%%%%%%%%%%%%%%%%%%%%%%%%%%%%%%%%%%%%%%%%%%%%%%%%%%%%%%%%%%%%%%%%%%%%%%%%%%

\paragraph{From Welfare to Change.}
Chapter~10 establishes what \emph{could} generate utility. But for behavioral 
change, we need:

\begin{itemize}
\item \textbf{Chapter 11 (Awareness):} What is the person \emph{actually thinking about}?
\item \textbf{Chapter 12 (Willingness):} Is the person \emph{ready to change}?
\end{itemize}

\paragraph{The Link.}
The welfare function $Q$ provides the foundation:

\begin{equation}
WTC = \sum_i a_i \cdot \omega_i \cdot \Delta U^i
\end{equation}

where $\Delta U^i = U^i_{new} - U^i_{old}$ comes from the welfare function, 
and $a_i$ (awareness) determines which components are salient in the decision moment.

%%%%%%%%%%%%%%%%%%%%%%%%%%%%%%%%%%%%%%%%%%%%%%%%%%%%%%%%%%%%%%%%%%%%%%%%%%%%%%%
\subsubsection{Chapter 10 Summary: The Complete FEPSDE-Utility Framework}
%%%%%%%%%%%%%%%%%%%%%%%%%%%%%%%%%%%%%%%%%%%%%%%%%%%%%%%%%%%%%%%%%%%%%%%%%%%%%%%

\begin{tcolorbox}[colback=gray!5!white, colframe=gray!75!black, title=Chapter 10 Complete Summary]

\textbf{The 144-Component Structure:}
$$144 = 3 \text{ (INU/KNU/IDN)} \times 6 \text{ (FEPSDE)} \times 4 \text{ (time)} \times 2 \text{ (Gain/Pain)}$$

\textbf{The Master Formula (Component Level):}
$$U^i_{d,t}(\Psi) = \delta^{i,t}_d(\Psi) \cdot \left[ G^i_{d,t} - \lambda^i_d(\Psi) \cdot P^i_{d,t} \right]$$

\textbf{The Aggregate Welfare Function:}
$$Q(\Psi) = \sum_{i} \omega_i \cdot U^i(\Psi) + \sum_{i < j} \gamma_{ij}(\Psi) \cdot U^i(\Psi) \cdot U^j(\Psi)$$

\textbf{Key Innovations:}
\begin{enumerate}
\item \textbf{Three utility categories} (INU, KNU, IDN) with identical structure
\item \textbf{Six welfare dimensions} (FEPSDE) capturing full human welfare
\item \textbf{Reference-dependent} evaluation with category-specific reference points
\item \textbf{Inter-category complementarities} determining utility coherence
\item \textbf{Context modulation} via $\Gamma$-matrix (48 coefficients)
\item \textbf{Structural trade-offs} revealed ($\gamma_{S,M} < 0$, $\gamma_{Eco,M} < 0$, $\gamma_{Eco,E} < 0$)
\end{enumerate}

\textbf{The K-Q-$\Psi$ Relationship:}
$$\Psi \longrightarrow K \longrightarrow Q$$
Context enables coherence; coherence enables welfare.

\textbf{Policy Principles:}
\begin{enumerate}
\item Optimize multi-dimensionally (not just GDP)
\item Balance utility categories (not just INU)
\item Calibrate to context (not one-size-fits-all)
\item Manage reference points (not just outcomes)
\item Sequence correctly: $\Psi$ first, then $\gamma$, then $x$
\end{enumerate}

\textbf{Foundation for EBF:}
\begin{itemize}
\item Chapter 11: Awareness (which components are salient?)
\item Chapter 12: Willingness ($WTC = \sum a_i \omega_i \Delta U^i$)
\item Chapter 13: Journey (how does change unfold?)
\item Chapter 14: Segments (who are the people?)
\end{itemize}

\end{tcolorbox}

%%%%%%%%%%%%%%%%%%%%%%%%%%%%%%%%%%%%%%%%%%%%%%%%%%%%%%%%%%%%%%%%%%%%%%%%%%%%%%%
% END OF CHAPTER 10
%%%%%%%%%%%%%%%%%%%%%%%%%%%%%%%%%%%%%%%%%%%%%%%%%%%%%%%%%%%%%%%%%%%%%%%%%%%%%%%
